% \iffalse meta-comment
%
% hit-hyperlink.dtx 是各个类共用的 hyperref 配置
%
% \fi
%
% \section{共享超链接配置}
%
% \changes{v3.2a}{2026/08/31}{hyperref 的配置抽出来共用；题目、作者、关键词写进 PDF 属性；
%   \cs{enskip} 在书签里换成空格（英文章名每份都报 Token not allowed）；钉住
%   \cs{Urlmuskip}}
%
% PDF 书签、链接边框、url 字体这些跟学校规范无关，是模板整体的取舍，三个类原先
% 各写一遍且一字不差。
%
% \cs{hit@setup@hyperlink} 只做配置，由各个类在原来那段代码所在的位置调用，
% 所以生成物里的加载顺序不变。\pkg{natbib} 的引用样式补丁没有一起搬——那是
% \file{hit-thesis.cls} 独有的，也不属于超链接。
%
%    \begin{macrocode}
%<*shared-hyperlink>
%    \end{macrocode}
%
% \begin{macro}{\hit@setup@hyperlink}
% \texttt{colorlinks=false} 加 \texttt{pdfborder=0 0 0}：屏幕上不上色也不画框，
% 打印出来与正文一致。\texttt{CJKbookmarks} 让中文书签不乱码，
% \texttt{bookmarksopenlevel=1} 只展开到章。
%    \begin{macrocode}
\cs_new_protected:Npn \hit@setup@hyperlink
  {
\RequirePackage{hyperref}
\hypersetup{%
  CJKbookmarks=true,
  linktoc=all,
  bookmarksnumbered=true,
  bookmarksopen=true,
  bookmarksopenlevel=1,
  breaklinks=true,
  colorlinks=false,
  plainpages=false,
  pdfborder=0~0~0
}
\urlstyle{same}
%    \end{macrocode}
% 链接内部的可伸量各家给的不一样：\pkg{url} 本体 0mu，\pkg{gbt7714} 各版
% 又搁自己的（老版 0mu~plus~0.1mu、新版 0mu~plus~3mu，还都挂在文档开头的
% 钩子里），实测 2026 文档里落定的是 0mu。伸量一动，同一条长网址的断点
% 就差一位，连带掀整段断行。压轴（begindocument/end）统一钉成 0mu。
% 断点集也一样：老版 \pkg{gbt7714} 往 \cs{UrlBreaks} 里塞进全部字母
% 数字（网址随处可断），新版收回标准集，同一条网址断的位置就不一样。
% 老版压轴时把 \cs{UrlBreaks} 回置成 \pkg{url} 的默认集。
%    \begin{macrocode}
\AddToHook { begindocument / end }
  {
    \Urlmuskip = 0 mu \scan_stop:
    \cs_if_exist:NF \__xeCJK_use_ecglue_skip:
      {
        \cs_set_nopar:Npn \UrlBreaks
          {
            \do \. \do \@ \do \\ \do \/ \do \! \do \_ \do \|
            \do \; \do \> \do \] \do \) \do \, \do \? \do \&
            \do \' \do + \do \= \do \#
          }
      }
  }
  }
%    \end{macrocode}
% \end{macro}
%
% \begin{macro}{\hit@setup@pdf@metadata}
% 元信息什么时候齐要看用户怎么写：示例里 \cs{input}\marg{front/cover} 在
% \cs{begin}\marg{document} 之后，所以挂 \cs{AtEndPreamble} 太早，字段还是空的。
% 挂在 \cs{makecover} 上最稳：封面要用的正是这几个字段，跑到那里一定齐了。
% 取的是 \cs{hit@title@cover@zh} 而不是 \cs{hit@title@flat@zh}：
% 后者为了屏蔽 \cs{\textbackslash\textbackslash} 包了一层 \cs{begingroup}，而
% \cs{begingroup} 进不了 PDF 字符串，\pkg{hyperref} 会逐个报警告。
%
% \cs{pdfstringdefDisableCommands} 是 \pkg{hyperref} 给的钩子，在把记号转成 PDF
% 字符串时生效，把 \cs{\textbackslash\textbackslash} 与 \cs{enskip} 换成空格。
% 后者出现在英文版的章名里（\option{chapter-name} 取 \texttt{Chapter\string\enskip,}），
% 不声明的话 \pkg{hyperref} 照样会换成空格，但要报一条警告。关键词取的是
% \cs{hit@keywords@zh@plain}——那一份没有 \cs{ignorespaces}，词之间只有分隔符，
% 与排出来的版面一致。关键词只有 \file{hit-thesis.cls} 有，所以先探一下。
%    \begin{macrocode}
\cs_new_protected:Npn \hit@setup@pdf@metadata
  {
    \pdfstringdefDisableCommands
      {
        \cs_set:Npn \\ { ~ }
        \cs_set:Npn \enskip { ~ }
      }
  }
\cs_new_protected:Npn \hit@write@pdf@metadata
  {
    \hypersetup
      {
        pdftitle  = { \hit@info@title@cover@zh } ,
        pdfauthor = { \hit@info@author@zh } ,
      }
    \cs_if_exist:cT { hit@info@keywords@zh }
      { \hypersetup { pdfkeywords = { \use:c { hit@info@keywords@zh@plain } } } }
  }
%    \end{macrocode}
% \end{macro}
%
%    \begin{macrocode}
%</shared-hyperlink>
%    \end{macrocode}
%
% \Finale
